Population-Weighted Compactness captures a dimension of district quality that
is orthogonal to geometric shape metrics.
Our principal findings are:

\begin{enumerate}
  \item PWC--PP correlation is $r < 0.4$ across all algorithm-state pairs,
    confirming that representational compactness (PWC) and geometric compactness
    (PP) measure independent properties. A complete compactness analysis requires both.
  \item The prime-factor structure algorithm reduces mean PWC approximately 14\%
    below standard-bisect on NC 2020 due to hierarchical nesting that aligns
    population concentrations with district centroids.
  \item TX exhibits much higher PWC variance than NC or WI due to the geographic
    concentration of population in major urban centres (Dallas-Fort Worth, Houston,
    San Antonio-Austin corridor), making PWC a state-specific rather than universal
    compactness benchmark.
\end{enumerate}

Practitioners should include PWC as the representational dimension in composite
compactness profiles (K.7), using the prime-factor algorithm when PWC minimisation
is a priority.
PWC is particularly valuable for VRA district analysis, where the clustering
of minority voters near a district centroid is a relevant indicator of
community cohesion.
